(1) 定义 \(q=x/\rho,\ p=\rho\dot x-\dot\rho x\)，则
\[
\dot q=\frac p{\rho^2},\qquad
\dot p=\rho\ddot x-\ddot\rho x
=-\rho\omega^2x-(\rho^{-3}-\rho\omega^2)x=-\frac q{\rho^2}.
\]
因此 \(\dot I=q\dot q+p\dot p=0\)。这没有使用绝热假设；只要辅助方程和原始方程成立，二次量便精确守恒。若外力调制频率，一般的机械能并不守恒，因而两者不可混为一谈。

(2) 由 \(\tau(t)=\int_0^t\rho(s)^{-2}ds\)，有 \(dq/d\tau=p,\ dp/d\tau=-q\)，故
\(q=A\cos\tau+B\sin\tau\)、\(x=\rho(A\cos\tau+B\sin\tau)\)。对应基本解
\(u=\rho\cos\tau,\ v=\rho\sin\tau\) 自动满足原运动方程，且
\[
W=u\dot v-\dot uv=\rho^2\dot\tau=1.
\]
它们线性独立，张成全部解空间，\(I=(A^2+B^2)/2\)。此时间重参数化把任意正频率允许的Pinney配对转化成单位频率振子。

(3) 写 \(f=\rho^2=a^2\cos^2\omega_0t+b^2\sin^2\omega_0t\)。直接求导得
\(\ddot\rho+\omega_0^2\rho=(ab\omega_0)^2/\rho^3=1/\rho^3\)。
在正切函数的每一个连续分支上，
\[
\tau(t)=\operatorname{unwrap}\arctan[(b/a)\tan(\omega_0t)]
\]
且以 \(\tau(0)=0\) 并跨极点连续延拓；因为 \(ab\omega_0=1\)，其导数等于 \(1/\rho^2\)。初值给 \(A=x_0/a,\ B=av_0\)，因而
\[
x(t)=\rho(t)[(x_0/a)\cos\tau(t)+av_0\sin\tau(t)].
\]
同一初值问题唯一性说明它与通常的 \(x_0\cos\omega_0t+(v_0/\omega_0)\sin\omega_0t\) 恒等。

(4) 对 \(H\) 微分并代入运动方程，
\[
\dot H=\dot x(\ddot x+\omega^2x)+\omega\dot\omega x^2
=\omega\dot\omega x^2.
\]
映射矩阵为
\(\begin{pmatrix}1/\rho&0\\-\dot\rho&\rho\end{pmatrix}\)，行列式恰为1，保持 \(dq\wedge dp=dx\wedge d\dot x\)。即使在恒频情况下，非平凡辅助解 \(a\ne b\) 的 \(I=(x_0/a)^2/2+(av_0)^2/2\) 也一般不等于机械能。若 \(a=b=1/\sqrt{\omega_0}\)，则 \(\rho\) 常数、\(I=H/\omega_0\)，成为自然的作用量归一化。
